%%% raster.tex — welche Elemente verlassen das Grundlinienraster?
%%%
%%% Jedes Element steht zwischen zwei Absaetzen Fliesstext. Wenn die Zeilen
%%% danach nicht mehr auf dem Raster sitzen, liegt es an dem Element
%%% dazwischen — und werkzeuge/nachmessen.py sagt, um wie viel.
%%%
%%%   xelatex raster.tex
%%%   python3 ../werkzeuge/nachmessen.py raster.pdf --seite 1 --text

\documentclass[raster,ohnehintergrund]{dsa5latex}
\graphicspath{{../}}

\begin{document}
\onecolumn

Erste Zeile vor allen Elementen, sie setzt den Bezug.
Zweite Zeile.
Dritte Zeile.

\begin{dsaliste}
\item Erster Punkt einer Aufzaehlung
\item Zweiter Punkt
\end{dsaliste}

Nach der Aufzaehlung, erste Zeile.
Nach der Aufzaehlung, zweite Zeile.

\begin{dsalisteblau}
\item Blaue Raute, erster Punkt
\end{dsalisteblau}

Nach der blauen Aufzaehlung.

\begin{dsaWerteabsatz}
Werteabsatz mit haengendem Einzug, erste Zeile, lang genug fuer einen Umbruch
in eine zweite Zeile.
\end{dsaWerteabsatz}

Nach dem Werteabsatz.

\dsaVorlesetext[9]{Vorlesetext mit der Zierleiste darueber.}

Nach dem Vorlesetext, erste Zeile.
Nach dem Vorlesetext, zweite Zeile.

\dsaEinfuehrung{Einfuehrungstext, kursiv mit Laufweite.}

Nach der Einfuehrung.

\dsaZitat{Freistehendes Zitat.}

Nach dem Zitat.

\dsaabschnitt{Abschnitt, 13 bp fett}

Nach dem Abschnitt.

\dsaunterabschnitt{Unterabschnitt, 10 bp fett}

Nach dem Unterabschnitt.

\dsaunterkapitel{Unterkapitel, 14 bp zentriert}

Nach dem Unterkapitel.

\dsaTabellenkopf{Ein Tabellenkopf}

\dsaTabellenlinie
\begin{dsaTabelle}{@{}ll@{}}
1 & erster Wert\\
2 & zweiter Wert\\
3 & dritter Wert\\
4 & vierter Wert\\
5 & fuenfter Wert\\
6 & sechster Wert\\
\end{dsaTabelle}

Nach der Tabelle.

\dsaProbe{Sinnesschaerfe (Suchen), erschwert um 2}

Nach der Probenzeile.

\end{document}
